\chapter{Rationality and the Value-Freedom of Science}
\label{sec:p11-rationality-14-2}

Values can affect which questions scientists investigate, which alternatives
they develop, how they evaluate evidence, and what they recommend doing.
These are different influences. Anke Bueter's chapter~14.2 asks whether
excluding non-epistemic values from theory assessment is a defensible ideal
of scientific rationality. Its answer is not that interests make any
conclusion acceptable. Rather, the separation of a value-free epistemic
core from the surrounding purposes and consequences of inquiry requires
arguments that the chapter examines and contests
\cite[pp.~757--765]{Bueter2021ValueFreedom}.

The source is a philosophical survey and argument, not an experiment on
scientists or a formal proof of an ethical conclusion. Its historical cases,
reviewed empirical examples, proposed norms and objections have different
standing. The discussion below retains those distinctions and the source's
eight notes. Its conclusion about an emerging rejection of the traditional
ideal is Bueter's 2021 assessment of a debate, not a measured consensus or
a verdict settling every account of science and values.

\section{Two directions of concern, and a restricted ideal}
\label{sec:p11-rat142-ideal}

\begin{hbdefinition}{def:p11-rat142-value}{Values, value-ladenness and the disputed norm}
Bueter uses value broadly for something considered desirable or worthy of
pursuit. Religious commitments, social ideals, economic interests,
political goals and individual preferences can all fall under that usage.
A value judgment need not be an arbitrary preference. Some proponents of
value-freedom allow rational discussion of values while denying that
scientific methods alone justify them.

A claim is value-laden in the sense at issue when its formulation,
evidential interpretation or methodological standards depend on evaluative
commitments. That description does not yet decide whether the dependence
is legitimate. The contemporary value-free ideal addressed here permits
values to influence the direction and application of research and the
ethical permissibility of methods, but requires epistemic assessment of
theories to exclude non-epistemic influences.
\index{value-freedom!science}
\index{science!epistemic and non-epistemic values}
\index{inductive risk}
\end{hbdefinition}

Bueter distinguishes an \emph{is-to-ought} question from an
\emph{ought-to-is} question. The first concerns whether empirical findings
justify political or ethical prescriptions. The second concerns whether
values influence judgments presented as empirical knowledge. Confusing
them obscures the argument: recognizing a legitimate role for values in
research does not make a desired outcome evidence that a theory is true;
denying an inference from facts alone to a prescription does not make
research priorities or evidential standards independent of values.

\subsection{Science, citizenship and political neutrality}

In the historical route presented by Bueter, Hume criticizes an unexplained
transition from descriptive assertions to obligations. Weber distinguishes
the scientific role from the citizen's political role. Science can study
value systems and their coherence, investigate whether a goal is achievable,
and assess available means. On Weber's account it cannot confer scientific
authority on the ultimate end itself. Bueter's example contrasts increasing
production with ensuring fair working conditions: economics may illuminate
efficient means to greater production without deciding that production
should take priority over fairness. The priority requires an additional
evaluative commitment, not merely a production forecast
\cite[sec.~1.1, pp.~757--758]{Bueter2021ValueFreedom}.

The 1960s positivism dispute supplies an objection and a reply. Critical
rationalists associated with Popper and Albert objected that the Frankfurt
school's political prescriptions borrowed scientific authority without
being exposed to the relevant empirical tests. Adorno, Habermas and others
argued that refusing social critique under the banner of neutrality can
preserve prevailing power relations; they regarded critique as a task of
social science. These are opposed claims about method, authority and social
responsibility. Neither the authority of a speaker nor the fact that a
position benefits an existing institution resolves them. The distinction
between explanation and normative warrant in
\cref{def:p11-reasons-roles,sec:p11-method-boundary} remains applicable.

\subsection{Discovery, justification and application}

The later debate focuses especially on influences upon scientific contents
and processes. Bueter reports feminist historical and philosophical work
showing gender prejudices in accepted theories. Such cases do not by
themselves refute value-freedom as a norm: a defender can regard them as
failures showing why the ideal should be pursued more vigorously.
An ideal need not be fully realized to remain a defensible goal. Bueter
uses Shrader-Frechette's defense of an unattained ideal and Koertge's
insistence on resisting political and religious influences in justification
to make that reply explicit. The chapter therefore turns from
counterexamples to systematic objections
\cite[sec.~1.2, pp.~758--759]{Bueter2021ValueFreedom}.

The familiar context distinction assigns hypothesis generation to discovery,
evidential testing to justification, and practical implementation to
application. Under the restricted ideal, a hypothesis's origin in a
misogynist commitment, extended reflection or a nightmare about snakes
does not determine its evidential
status. This blocks a simple genetic dismissal: one must evaluate the
claim and its support rather than condemn it solely by its origin.
Application can be delegated to democratic deliberation and policy-makers,
while ethical restrictions can govern admissible experimental methods.

This defense depends on more than naming three stages. It assumes that
the justified outcome can be adequately assessed independently of the
values determining what is developed, measured or subsequently done.
Bueter's note~3 uses the distinction to organize the controversy, not to
endorse that independence or a rigid partition of actual research. Her
note~2 similarly warns that the contemporary ideal has historical variants.
The target is not the implausible thesis that no scientist ever has a
practical interest, but a specific claim about legitimate theory appraisal.

\section{Discovery, pursuitworthiness and significant truth}
\label{sec:p11-rat142-pursuit}

The source's first systematic challenge concerns the range of scientific
rationality. Even without a unique algorithm for discovering hypotheses,
the choice of questions and the development of candidates need not be
irrational. They can be assessed in light of expected epistemic and
practical prospects. The unrestricted aim of truth does not specify which
of the indefinitely many truths deserve scarce research resources.
Kitcher's account instead emphasizes \emph{significant} truths,
predictions and explanations. Explaining the solar system or forecasting
climate change involves a particular aim, not simply acquiring any true
sentence. What is treated as epistemically significant can itself reflect
the history of inquiry and earlier value-laden priorities
\cite[sec.~2, p.~759]{Bueter2021ValueFreedom}.

Kitcher's ideal of well-ordered science makes research direction publicly
accountable through deliberation representing different social interests,
informed by impartial scientific expertise. Bueter presents this as a
normative institutional proposal. It does not entail that an actual vote
determines whether a scientific proposition is true, or that every
deliberative institution successfully represents all affected interests.
The proposal broadens the assessment of rationality beyond final theory
acceptance without replacing evidential justification by popularity.

Getting from a promising question to a testable theory requires choices
about hypothesis formulation, variables, concepts, empirical consequences
to test and methods of testing them. It also requires funding, time and
the ability to continue after setbacks. A \emph{context of pursuit} makes
those choices visible: a candidate may be worth developing because of
plausibility, fertility or expected epistemic/practical promise before
it is justified as true. Laudan's distinction between pursuing and accepting
a theory therefore concerns two objects of appraisal. Funding a promising
candidate is not equivalent to endorsing its present evidential adequacy.

The second challenge is that these earlier choices affect later assessment.
Okruhlik's example supposes a thoroughly sexist setting in which nonsexist
rivals are not developed. Applying rigorous comparative standards among
the generated theories can then select the best \emph{available} sexist
rival without removing the agenda's earlier influence. Elliott and
McKaughan's reviewed toxicology cases make a related point about which
theories and data become available. Bueter's blind-spot argument adds that
an ungenerated alternative or uncollected observation might otherwise have
undermined the accepted conclusion. The premises concern the actual range
of candidates and evidence; they do not say that a merely imagined rival
already defeats every established result
\cite[pp.~759--760]{Bueter2021ValueFreedom}.

\begin{hbexample}{def:p11-rat142-available-rivals}{Rigorous comparison on a restricted agenda}
Suppose a research program compares two drug-based treatments using sound
methods, but no one investigates a feasible non-drug alternative because
funding priorities systematically exclude it. A conclusion that one tested
treatment performed better than the other can be warranted for that
comparison. The broader conclusion that it is the best available kind of
treatment has not been established: the unstudied option is absent from
the comparison, and might be better, worse or equivalent.

This is a qualitative reconstruction of the chapter's omitted-alternative
mechanism, not a report of a clinical trial or a recommendation about a
particular treatment. A truthful restricted comparison can become misleading
when communicated with a broader scope. It also illustrates Anderson's
point, as presented by Bueter, that a partial truth can mislead without
being a literally false result. Better testing of the same two options
does not itself repair the agenda omission.
\end{hbexample}

The chapter mentions profitable pharmaceutical research that neglects
nonpharmaceutical alternatives and cites Jukola's work on commercial
agenda-setting. This is a reviewed illustration, not evidence here that
every commercially funded study is unreliable or that an untested alternative
works. To establish a particular commercial influence one would need its
research choices, available alternatives, evidence and competing explanations.
An absence can have many causes besides deliberate suppression.

This argument need not make the origin of every hypothesis decisive.
A distinction remains between using a hypothesis's origin as a substitute
for evaluating its evidence and examining how the process generating
candidates changes the content of the comparative test. Bueter's note~4
also identifies semantic value-ladenness of scientific concepts,
classifications and terms as another debate, but expressly does not develop
it. The present agenda analysis is not a complete account of that separate
issue.

\section{Epistemic values and contested boundaries}
\label{sec:p11-rat142-values}

The value-free ideal does not normally exclude every value from justification.
Its central distinction is between epistemic or cognitive criteria and
non-epistemic commitments. Bueter uses epistemic consistently for readability,
as her note~5 explains, while retaining the substantive disagreement about
what the category comprises.

\begin{hbdefinition}{def:p11-rat142-kuhn-values}{Kuhn's theory-choice values in the source}
Kuhn's five criteria are accuracy or empirical adequacy; consistency,
both internal logical consistency and compatibility with other accepted
theories; scope; simplicity; and fruitfulness. They guide the appraisal
of theories rather than supply an algorithm selecting one in every case.
Their interpretation and relative importance can differ with the research
context and the individuals making the assessment. Calling them values
does not make empirical adequacy optional or identify them with political
preferences.
\end{hbdefinition}

McMullin calls such criteria epistemic because of their contribution to
truth-like science and contrasts them with religious or sociopolitical
values. Rooney challenges the clarity of that boundary: different lists
of epistemic values may incorporate social commitments into what counts
as epistemic success. Bueter does not treat the existence of multiple
lists as a mathematical proof that no useful distinction can be drawn;
it motivates examining a criterion's function in a particular context
\cite[sec.~3, pp.~760--761]{Bueter2021ValueFreedom}.

Longino's alternative begins from the goals of inquiry. Research intended
to understand gender oppression or develop nonsexist accounts of women's
health can have aims that are both feminist and epistemic. Alongside
internal consistency and empirical adequacy, her alternative list adds
novelty, ontological heterogeneity, complexity of relation, applicability
to human needs and diffusion of power. In some contexts the Kuhnian
criterion of simplicity can favor an androcentric theory, for instance
when simplification excludes women's perspectives or concerns.
The contextual claim is not that simplicity always causes sexism or
that adding complexity makes a theory true.

Note~7 is important to the comparison. Neither Kuhn nor Longino is here
committed to every rival being strictly empirically equivalent. Kuhn's
moderate underdetermination allows different candidates to fit different
areas of evidence better; Longino emphasizes that substantive and
methodological background assumptions mediate evidential relations and
may be value-laden. Agreement on some minimum criteria can coexist with
disagreement on how to choose among candidates. It does not imply that
evidence places no restrictions on choice
\cite[pp.~760,763, notes~6--7]{Bueter2021ValueFreedom}.

The chapter preserves defenses of a more discriminating epistemic category.
Ruphy proposes evaluating epistemic values by their contribution to
empirically successful theories. Douglas distinguishes two axes: ideal
desiderata versus minimal criteria, and properties of a theory itself
versus properties of the theory in relation to evidence. Some resulting
groups are presented as directly truth-assuring; others are strategic or
pragmatic, making defects more likely to be found through such features
as testability. Thus not all epistemic values need have exactly the same
connection to truth.

Lacey separates three attitudes often compressed into accepting a theory:
adopting it to direct further research; endorsing it as sufficiently
supported for practical use; and impartially holding it of specified
phenomena as established without a need for further testing. He confines
legitimate influences in the third attitude to epistemic values. Rooney
objects that this presupposes the epistemic/non-epistemic demarcation
instead of grounding it. The distinction between attitudes clarifies
what is being evaluated; it does not by itself justify the proposed
restriction on that evaluation.

Steel differentiates intrinsic epistemic values, necessary for truth
attainment, from extrinsic ones promoting it indirectly. Bueter's examples
are empirical adequacy in the first group and testability or the ability
to obtain funding in the second. Calling a funding capacity extrinsically
epistemic concerns its enabling role in inquiry, not a general claim that
profit incentives are truth-promoting. Under the reported account a
non-epistemic value is irrelevant or harmful to truth attainment. These
are one author's functional distinctions, not uncontroversial classifications
of every social aim.

\subsection{Epistemic constraints versus inquiry's particular aims}

Even a defensible classification leaves a normative question about priority.
An epistemic-constraints approach, represented by Steel, prohibits
non-epistemic values from overriding adequate-science requirements while
permitting them as tie-breakers when epistemic criteria leave a choice
underdetermined. An aim-oriented approach, represented by Elliott and
McKaughan, starts from inquiry's particular goals, integrating significant
truths with practical concerns rather than granting an unqualified priority
to an abstract aim of truth. Values contributing to those goals can have
legitimate roles if goals and influences are made maximally transparent.
Brown argues that avoiding wishful thinking does not require blanket
epistemic priority when value judgments themselves are open to rational
criticism and revision
\cite[p.~761]{Bueter2021ValueFreedom}.

These proposals impose different conditions, not different names for an
identical rule. A social aim used only to choose between two scientifically
adequate candidates fits the tie-breaker proposal. A proposal allowing that
aim to shape which epistemic achievements matter goes further and requires
a defense of the integrated goal. Transparency identifies the commitment;
it does not prove that the commitment is legitimate or that the reported
evidence supports the favored answer. Bueter leaves open how to legitimate
epistemic-social goals and protect inquiry against corruption by political
or economic interests. The positive critique of a rigid ideal is therefore
not a license for sponsor preference to settle a fact.

\section{Inductive risk and the consequences of error}
\label{sec:p11-rat142-inductive-risk}

\begin{hbdefinition}{def:p11-rat142-risk-profile}{Hypothesis acceptance, errors and evidential standards}
Write \(H\) for the descriptive hypothesis under assessment. In the
acceptance terminology used by Bueter's account of Rudner, a false positive
accepts \(H\) when it is false; a false negative rejects \(H\) when it is
true. Assessing the consequences of those errors and deciding how much
evidence suffices are separate from deciding what the evidence actually is.

For comparison with \cref{def:p15-statistical-experiment},
\(\alpha\) denotes a selected test level: an upper bound on rejecting
the specified null under its sampling models. Mapping that statistical
type-I error to a false acceptance in the present case requires specifying
which proposition is null and which inference follows from rejection.
The word positive alone does not fix that mapping. Nonrejection of a null
is not proof of its truth or automatic acceptance of the opposite claim.
\end{hbdefinition}

Rudner's argument, as reconstructed by Bueter, proceeds through four
commitments. Scientists accept or reject hypotheses under uncertainty;
both kinds of error are possible; the decision requires a standard for
how much evidence is enough; and the significance of errors for people
affected by subsequent use is relevant to that standard. The last step
introduces evaluative considerations about acceptable risk. It is not
derived merely from the existence of an error probability. Given those
premises, Rudner concludes that scientists make value judgments in their
scientific role. A critic can contest whether acceptance is required,
whether it is an action-like judgment, who should set the standard, or
which consequences are foreseeable
\cite[sec.~4, pp.~761--762]{Bueter2021ValueFreedom}.

\begin{hbexample}{def:p11-rat142-drug}{The source's new-drug safety case}
Let \(H\) say that a new drug is safe under a specified descriptive
criterion. Suppose acceptance will probably lead to approval and widespread
use. Accepting a false safety hypothesis may expose many patients to severe
adverse effects. Rejecting a true one may forgo an effective treatment and
the company's profits. Deciding which error matters more invokes judgments
about harms, benefits, interests and affected groups.

A high evidential demand motivated by avoiding false safety assurances
uses values to assess how uncertainty is tolerated. Treating the desirability
of an inexpensive treatment as evidence of safety instead makes a
wished-for conclusion do evidential work it cannot do. Choosing a low
demand because approval is commercially attractive can still operate
indirectly through a standard, yet be normatively problematic. The difference
between direct and indirect roles therefore does not settle the justification
of every threshold choice.

The source supplies a conceptual illustration, not a named drug, clinical
dataset, estimated adverse-event rate or regulatory rule. Safety's descriptive
criterion and the subsequent approval policy must themselves be stated.
No medical recommendation or universal evidential threshold follows from
the example. Forgone treatment can be serious, so the error asymmetry is
not settled merely by labelling one side positive.
\end{hbexample}

Levi objects that Rudner blurs epistemically accepting a hypothesis with
acting on it. Jeffrey adds that scientists often cannot predict the eventual
uses of their work; unknown consequences cannot straightforwardly determine
a particular acceptance policy. He proposes reporting probabilities rather
than accepting or rejecting outright, leaving users to decide what is
enough for their circumstances. Betz similarly advocates hedged hypotheses
that express the amount and kind of uncertainty and objects to scientists
displacing democratic decisions about acceptable risk. Bueter also points
to John's reply to that defense. Reporting a probability, stating a hedged
claim, accepting a belief and recommending action are thus four options
in the controversy, not interchangeable descriptions of one task.

The distinctions in \cref{def:p11-rat22-attitude-types} make the first
objection precise: evidence bearing on a belief and reasons for acting
on it may concern the same proposition but evaluate different attitudes.
The posterior construction in \cref{sec:p15-bayesian-updating-prediction}
likewise supplies conditional probabilities given a model and prior,
not a rule assigning legitimate losses or requiring outright scientific
acceptance. Neither interface proves Jeffrey's philosophical proposal,
but both expose what additional commitments a claim requires.

\subsection{Douglas's strengthened argument and its ethical premise}

Douglas's account limits the application argument to policy-relevant
research in which important consequences of assertions can often be
foreseen. It also relocates risk-sensitive decisions earlier than final
acceptance: selecting significance standards, gathering and interpreting
data, and interpreting outcomes can change the balance of error
probabilities. Deferring the final policy decision does not automatically
externalize every such choice. Her normative conclusion is that scientists
have an ethical obligation to consider potential practical consequences
of error when setting evidential standards. Foreseeable consequences and
that professional-responsibility premise matter to the inference. Their
mere foreseeability does not mathematically prove an obligation
\cite[p.~762]{Bueter2021ValueFreedom}.

\begin{hbdefinition}{def:p11-rat142-value-roles}{Direct and indirect roles in Douglas's proposal}
A direct role treats a value as a reason to accept or reject the claim
itself. An indirect role uses a value to assess how important the remaining
uncertainty is and hence what should count as sufficient evidence.
Douglas's proposed norm permits the latter role and rejects the former
when non-epistemic desirability substitutes for evidential support.
This is a proposed distinction about how a value operates, not a claim
that one can identify legitimacy by looking only at its subject matter.
\end{hbdefinition}

Bueter records objections to both conceptual sharpness and normative
sufficiency of this distinction. An indirect influence can be problematic
when it biases the error balance toward a sponsor's interests. Making a
favored result easier to declare by adjusting a methodological threshold
does not become legitimate merely because the alteration is described
as indirect. Conversely, explaining a risk-sensitive threshold openly
does not by itself prove that the chosen risks are acceptable. Different
effects on different groups, the available evidence and the aims of the
research still require assessment.

The source reports debate over the inductive-risk argument's scope and
soundness and mentions Biddle and Kukla's call for finer distinctions among
epistemic risks. Examples ranging from medicine to climate science and
theoretical physics demonstrate the breadth of the literature reviewed;
they do not establish one invariant error cost or that a single acceptance
policy fits all those domains.

\section{Conventions, social objectivity and public trust}
\label{sec:p11-rat142-objectivity}

Wilholt's alternative, as presented by Bueter, emphasizes methodological
conventions shared by a scientific community. They standardize decisions
that can affect acceptance indirectly. Departures may make value influences
visible and matter to epistemic trust. A deviation is an occasion for
examining its justification, not a deductive proof of misconduct; an accepted
convention is not beyond revision merely because it is shared. The account
also extends the error-balance problem beyond foreseeable practical harm:
false positives and false negatives can differ in epistemic importance
because a hypothesis has different consequences for a system of beliefs
\cite[p.~762]{Bueter2021ValueFreedom}.

\begin{hbexample}{def:p11-rat142-indirect-bias}{The same indirect mechanism, different justification}
Suppose a field employs a common evidential convention. Two researchers
depart from it before evaluating a policy-relevant claim. One explains how
the conventional balance poorly represents a newly identified serious
error consequence and submits the alternative standard and its rationale
to scrutiny. The other adjusts the standard only so a sponsor-favored
claim clears it, suppressing the justification and sensitivity of the result.

Both changes operate through evidential standards and are indirect in
Douglas's sense. That description does not make their justification the
same. The first still needs defensible evidence, values and a discussion
of competing consequences; it is not certified by openness alone. The
second shows why indirectness is insufficient as a normative safeguard.
This is a hypothetical contrast illustrating the chapter's objection,
not a measurement of how often sponsors influence research.
\end{hbexample}

Bueter's note~1 distinguishes value-freedom from objectivity: it has
traditionally been proposed as necessary but not sufficient for the latter.
Removing a particular interest does not ensure sound sampling, accurate
measurement or correct inference. Conversely, note~6 introduces Longino's
social objectivity, which organizes a community for effective critical
discourse rather than presupposing the absence of values. The objectivity
treatment's two senses in \cref{def:p11-rat141-objectivity} distinguish
whether theories faithfully describe the world from whether methods are
independent of personal and social values. Its Longino-style critical
process in \cref{def:p11-rat141-critical-process} concerns effective
community criticism and response, not merely increasing the number of
perspectives. These interfaces locate Bueter's proposal without treating
value-freedom, true results and functioning critical institutions as the
same property. Her questions about agenda, epistemic classification and
error responsibility do not reduce to a general objectivity label.

Note~8 addresses an objection to pluralism about epistemic values: if
standards differ, how can dissent be discussed rationally? Carrier's
answer locates a common epistemic attitude in procedures for debating
knowledge claims and seeking agreement, rather than requiring identical
theory-choice standards. This is another institutional proposal. It does
not make every disagreement resolvable, treat consensus as sufficient
evidence, or establish that procedural agreement always legitimates the
research aims being pursued \cite[p.~763, notes~1,6--8]{Bueter2021ValueFreedom}.

The concluding public-trust problem asks whose values enter inquiry,
how those influences are legitimated, and how the resulting science can
deserve rational trust. McKaughan and Elliott advocate transparency;
John warns that increased disclosure can erode public trust. Bueter
retains this tension rather than concluding that disclosure invariably
increases or invariably decreases trust. Trustworthiness, observed trust
and a justified decision to trust are different objects. A communication
strategy's effects require empirical evidence about audiences, messages,
designs and measures; the chapter provides no population estimate or
controlled intervention effect for such a strategy
\cite[secs.~5, pp.~762--763]{Bueter2021ValueFreedom}.

\section{Using the source's distinctions in a research decision}
\label{sec:p11-rat142-use}

\begin{hbmethod}{def:p11-rat142-assessment}{A value-influence assessment with explicit boundaries}
First identify the judgment: choosing a question, pursuing a candidate,
estimating support, accepting a belief, endorsing practical use or making
a policy recommendation. Record its proposition, evidential base and scope.
An unspecified reference to accepting a theory leaves the central issue open.

Identify the operative goal and the values contributing to it. State which
are claimed to be epistemic, why, and under which classification. Do not
use the category as its own justification. For a comparative appraisal,
ask which alternatives were developed and which were left out, what data
were collected, and how those choices bear on the conclusion's scope.

Identify the uncertainty and the consequences of error for the relevant
people or belief system. Specify the null/alternative and decision meanings
before translating statistical error labels. Explain who set the evidential
standard, what convention it follows or departs from, and which responsibility
or legitimacy premise supports that choice.

Distinguish a value's direct substitution for evidence from its indirect
influence on standards; then examine the indirect role rather than declaring
it automatically legitimate. Make the relevant goals, assumptions, funding
interests and alternatives open to criticism. Transparency supports
inspection; it does not supply a missing evidential or normative premise.

Finally state separately what the evidence establishes, what remains uncertain,
what recommendation follows under the adopted values, and which premises
are contested. If the inquiry supplies only a restricted comparison, do
not communicate an unrestricted result. If values remain disputed, the
outcome is an explicit conditional or disagreement, not a theorem of
scientific irrationality.
\end{hbmethod}

This editorial assessment procedure operationalizes the chapter's distinctions
without attributing a universal algorithm to Bueter. A successful use makes
the production/fairness priority, omitted-treatment comparison or drug-error
standard assessable without converting a preference into a fact. Failure
cases include inferring superiority over untested alternatives, adopting
a desired safety claim without evidence, treating indirect sponsor influence
as automatically legitimate, or claiming that disclosed values alone ensure
objectivity and trust.

The formal test results of \cref{sec:p15-hypothesis-tests-confidence}
hold under stated sampling models, error constraints and alternatives.
They do not choose the ethically correct constraint. Similarly,
\cref{def:p17-information-policy} compares policies given a prior and
payoff scale; a mathematical optimum within that problem does not prove
the legitimacy of the payoff assignment or whose interests it represents.
The source's inductive-risk debate concerns those interpretation and
responsibility questions, not a refutation of conditional statistical or
decision theorems.

The wider science/value discussion in \cref{sec:p11-ri-frontiers}
separates research purposes from epistemic method. Bueter adds a specific
challenge to keeping those stages independent. Her conclusion is that
legitimate and illegitimate value influence must be distinguished with
more refined accounts of roles, stages, goals and institutions. The defense
of empirical authority against wishful thinking survives that challenge;
the claim that any non-epistemic influence makes science irrational does not
follow from it. Questions about legitimate goals, acceptable risks and
democratic authority remain philosophically contested.

The chapter reviews empirical and historical cases but supplies no original
participant study, clinical data table or numerical effect estimate. Those
cases support the author's philosophical discussion at the reported scope,
not new quantitative generalizations. Stronger claims about a particular
scientific community, sponsor, research program or public audience require
the task/population/design/measure/uncertainty and alternative-explanation
record of \cref{def:p19-evidence-record}. The selected source is read in
full; its further books and case studies remain distinct presentations,
not proof or evidence silently imported by a reference list.
